\title{Title of Introduction}
\label{sec:intro}
\maketitle

\section{Title of Section}
Lorem ipsum dolor sit amet, consectetur adipiscing elit \cite{doe2020lorem}.  
Sed do eiusmod tempor incididunt ut labore et dolore magna aliqua \cite{roe2018ipsum}.  
Ut enim ad minim veniam, quis nostrud exercitation ullamco laboris nisi ut aliquip ex ea commodo consequat.\footnote{This is a sample footnote.}  

\subsection{Equation Example}

A dummy equation to illustrate formatting:
\begin{equation}
    y = a x^2 + b x + c,
\end{equation}
where $a, b, c \in \mathbb{R}$ are arbitrary constants \cite{smith2015dolor}.

\subsection{Acronyms Example}
\begin{itemize}
    \item This is an acronym with default behaviour: \gls{ai}.
    \item This is short only: \acrshort{ml}.
    \item This is long only: \acrlong{ai}.
    \item This is full version: \acrfull{ml}.
\end{itemize}

\subsection{Theorem and Remark}

\begin{theorem}[Dummy theorem]
If $a, b > 0$, then $a+b > 0$.
\end{theorem}

\begin{remark}
This is a simple remark environment to highlight comments or explanations.
\end{remark}

\subsection{Lists Example}

Here is a bullet list:
\begin{itemize}
    \item First bullet
    \item Second bullet
    \item Third bullet
\end{itemize}

Here is a numbered list:
\begin{enumerate}
    \item First item
    \item Second item
    \item Third item
\end{enumerate}

\subsection{Table Example}

Here is a simple table (Table~\ref{tab:example}):

\begin{table}[h!]
\centering
\begin{tabular}{|c|c|c|}
\hline
Column 1 & Column 2 & Column 3 \\
\hline
A & B & C \\
1 & 2 & 3 \\
X & Y & Z \\
\hline
\end{tabular}
\caption{A dummy example table.}
\label{tab:example}
\end{table}

\subsection{Figure Example}

Here is a dummy figure (Figure~\ref{fig:example}):  

\begin{figure}[h!]
\centering
\includegraphics[width=0.4\textwidth]{example-image} 
\caption{A sample figure}
\label{fig:example}
\end{figure}

%% -------------------------------
%% CHAPTER-BASED BIBLIOGRAPHY
%% -------------------------------
%% Uncomment the lines below ONLY if you want a separate bibliography
%% at the end of THIS chapter (paper-based thesis style).
% \printbibliography[heading=subbibliography, title={References}]
% \stopcontents[chapters]
